\section{TikZ Extensions}

\sectionPage{TikZ Extensions}

\begin{frame}[c]{TikZ Extensions}{Overlay Effects in TikZ Images}
    \centering
    \begin{tikzpicture}[thick]
        \draw[visible on=<1->, onslide=<5->uhhred, -latex] (0, 0) -- (1, 0);
        \draw[visible on=<2->, onslide=<6->uhhgreen, -latex] (1, 0) -- (1, 1);
        \draw[visible on=<3->, onslide=<7->uhhblue, -latex] (1, 1) -- (0, 1);
        \draw[visible on=<4->, onslide=<8->uhhstone, -latex] (0, 1) -- (0, 0);
    \end{tikzpicture}
    \vspace{2ex}
    \begin{itemize}
        \item<1-> \texttt{visible on=<num-num>}: make node edge visible on that slide
        \item<5-> \texttt{onslide=\{<num-num>style\}}: apply style on given slides
    \end{itemize}
\end{frame}

\begin{frame}[c]{TikZ Extensions}{Standalone Legends}
    \centering
    \begin{tikzpicture}
        \begin{customlegend}[%
                legend entries={%
                    Entry 1,
                    Entry 2,
                    Entry 3 with comments,
                    Entry 4,
                },
                legend columns={2},
                legend cell align={left},
                legend style={draw=none},
            ]
            \addlegendimage{thick, uhhblue};
            \addlegendimage{area legend, thick, draw=none, fill=uhhred};
            \addlegendimage{surf, thick};
            \addlegendimage{dashed, uhhstone, mark=*}
        \end{customlegend}
    \end{tikzpicture}
    \vspace{2ex}
    \begin{itemize}
        \item Custom legends using \texttt{customlegend} environment
    \end{itemize}
\end{frame}

\begin{frame}[c]{TikZ Extensions}{A Sine in TikZ}
    \centering
    \begin{tikzpicture}
        \begin{axis}[
                scale only axis,
                axis x line=middle,
                axis y line=middle,
                xlabel={$t$},
                ylabel={$x_a(t) = A \cos(2\pi f - \theta)$},
                ymin=-2,
                ymax=2,
                width=6cm,
                height=4cm,
                xtick={0},
                ytick={1},
                yticklabels={$A$},
                yticklabel style={right},
            ]
            \addplot[blue, mark=none, domain=-4:4, samples=500, name path=cos] {cos(deg(2*x + 1))};

            \path[name path=yaxis] (axis cs:0, -2) -- (axis cs:0, 2);

            \path [name intersections={of=cos and yaxis,by=E}];

            \node[fill=blue, circle, inner sep=1pt, pin=30:$A\cos(\theta)$] at(E) {};

            \draw (axis cs:-2.07, -0.75) -- (axis cs:-2.07, -1.25);
            \draw (axis cs:1.07, -0.75) -- (axis cs:1.07, -1.25);
            \draw[latex-latex] (axis cs:-2.07, -1.25) -- node[below] {$T = 1/f$} (axis cs:1.07, -1.25);
        \end{axis}
    \end{tikzpicture}
    \vspace{2ex}
    \begin{itemize}
        \item Demo on how to plot a sine in TikZ
    \end{itemize}
\end{frame}

\begin{frame}[c]{TikZ Extensions}{An STFT Framework in TikZ}
    \centering
    \pgfplotsset{ax_stft/.style={%
	scale only axis,
	axis on top,
	xtick={\empty},
	ytick={\empty},
	xticklabels={\empty},
	yticklabels={\empty},
	width=0.3\textwidth,
	height=7ex,
	axis lines=middle,
	enlarge x limits=0.1,
	enlarge y limits=0.1,
	every axis x label/.style={at={(current axis.right of origin)},anchor=south west},
	every axis y label/.style={at={(current axis.above origin)},anchor=south}
	}
}%
%
\begin{tikzpicture}[
	box/.style={draw, fill=uhhblue, text=white, text width=5.5em, minimum height=4.5em, align={center}},
	frame annotation/.style={text=black, font=\footnotesize, above},
	frame boundaries/.style={fill=white, font=\footnotesize, inner sep=1pt}
]

\begin{axis}[
	ax_stft,
	xlabel={time},
	ylabel={noisy speech},
	name={noisy},
]
	\addplot[solid, uhhblue] table[x=sample, y=noisy] {data/stft.tsv};
\end{axis}

\begin{axis}[
	ax_stft,
	xlabel={time},
	at={($(noisy.south) - (0,1ex)$)},
	anchor=north,
	clip=false,
]
	\addplot[densely dotted, uhhblue] table[x=sample, y=win1] {data/stft.tsv} node[pos=0.25, frame annotation] {$\ell - 2$};
	\addplot[dashed, uhhblue] table[x=sample, y=win2] {data/stft.tsv} node[pos=0.5, frame annotation] {$\ell - 1$};
	\addplot[solid, uhhblue] table[x=sample, y=win3] {data/stft.tsv} node[pos=0.75, frame annotation] {$\ell$};

	\coordinate (noisy frame center) at(axis cs: -256, 0);
	\coordinate (noisy frame end) at(axis cs: -511, 0);
	\coordinate (noisy frame start) at(axis cs: 0, 0);
\end{axis}

\node[box, below=6ex of noisy frame center] (dft) {DFT};
\node[box, right=of dft, onslide=<2->{fill=uhhgreen}] (enhancement) {speech\\ enhancement};
\node[box, right=of enhancement] (idft) {IDFT};

\begin{axis}[
	ax_stft,
	xlabel={time},
	ylabel={enhanced speech},
	name={clean},
	at={($(idft.east |- noisy) - (2ex,0)$)},
	anchor=origin,
]
	\addplot[solid, uhhblue] table[x=sample, y=clean] {data/stft.tsv};
\end{axis}

\begin{axis}[
	ax_stft,
	xlabel={time},
	at={($(clean.south) - (0,1ex)$)},
	anchor=north,
	clip=false,
]
	\addplot[densely dotted, uhhblue] table[x=sample, y=win1] {data/stft.tsv} node[pos=0.25, frame annotation] {$\ell - 2$};
	\addplot[dashed, uhhblue] table[x=sample, y=win2] {data/stft.tsv} node[pos=0.5, frame annotation] {$\ell - 1$};
	\addplot[solid, uhhblue] table[x=sample, y=win3] {data/stft.tsv} node[pos=0.75, frame annotation] {$\ell$};

	\coordinate (clean frame center) at(axis cs: -256, 0);
	\coordinate (clean frame end) at(axis cs: -511, 0);
	\coordinate (clean frame start) at(axis cs: 0, 0);
\end{axis}

\draw[-] (noisy frame end) -- node[frame boundaries] {$K - 1$} (noisy frame end |- dft.north);
\draw[-] (noisy frame start) -- node[frame boundaries] {$0$} (noisy frame start |- dft.north);
\draw[-latex] (noisy frame center) -- (noisy frame center |- dft.north);

\draw[-] ($(dft.east) - (0, 4ex)$) -- node[frame boundaries] {$K - 1$} ($(enhancement.west) - (0, 4ex)$);
\draw[-] ($(dft.east) + (0, 4ex)$) -- node[frame boundaries] {$0$} ($(enhancement.west) + (0, 4ex)$);
\draw[-latex] (dft) -- node[yshift=1.2ex, font=\small] {$Y_{k, \ell}$} (enhancement);

\draw[-] ($(enhancement.east) - (0, 4ex)$) -- node[frame boundaries] {$K - 1$} ($(idft.west) - (0, 4ex)$);
\draw[-] ($(enhancement.east) + (0, 4ex)$) -- node[frame boundaries] {$0$} ($(idft.west) + (0, 4ex)$);
\draw[-latex] (enhancement) -- node[yshift=1.2ex, font=\small] {$\hat{S}_{k, \ell}$} (idft);

\draw[-] (idft.north -| clean frame end) -- node[frame boundaries] {$K - 1$} (clean frame end);
\draw[-] (idft.north -| clean frame start) -- node[frame boundaries] {$0$} (clean frame start);
\draw[-latex] (idft.north -| clean frame center) -- (clean frame center);

\node at($(noisy.north) + (-4em, 2em)$) (spk noisy) {\includegraphics[width=1.5em]{graphics/PersonRed.pdf}};
\node at($(noisy.north) + (0em, 2.5em)$) (noise) {\includegraphics[width=1.5em]{graphics/Monkey.pdf}};

\draw[red, decoration={expanding waves, angle=30, segment length=0.5ex}, decorate] (spk noisy) -- (noisy);
\draw[brown, decoration={expanding waves, angle=30, segment length=0.5ex}, decorate] (noise) -- (noisy);

\node at($(clean.north) + (-4em, 2em)$) (spk clean) {\includegraphics[width=1.5em]{graphics/PersonRed.pdf}};

\draw[red, decoration={expanding waves, angle=30, segment length=0.5ex}, decorate] (spk clean) -- (clean);

\end{tikzpicture}%
%

\end{frame}

\begin{frame}[c, fragile]{TikZ Extensions}{An STFT Framework in TikZ}
    \begin{itemize}
        \item Some notes on the STFT framework picture
        \item It requires the following files
            \begin{itemize}
                \item \verb|data/STFT.tex|
                \item \verb|graphics/PersonRed.pdf|
                \item \verb|graphics/Monkey.pdf|
            \end{itemize}
        \item Changing the folders of these files requires editing of the graphic itself
        \item The graphic itself is given in \verb|graphics/STFT.tex|
    \end{itemize}
\end{frame}

% \begin{frame}[c]{TikZ Extensions}{Image plots}
%     \centering
%     \begin{tikzpicture}[]
%         \begin{groupplot}[
%                 group style={%
%                     group size={2 by 2},
%                     xlabels at=edge bottom,
%                     ylabels at=edge left,
%                     horizontal sep=2.5cm,
%                 },
%                 axis on top,
%                 scale only axis,
%                 ymin=0,
%                 ymax=4,
%                 xmin=1.75,
%                 xmax=3.5,
%                 colormap name=viridis,
%                 xlabel={$t~/~\text{s}$},
%                 ylabel={$f~/~\text{kHz}$},
%                 width=0.3\linewidth,
%                 height=2.5cm,
%                 colorbar,
%                 colorbar style={width=2ex, xshift=-2ex},
%                 title style={anchor=south, at={(0.5, 0.9)}},
%             ]
%             \nextgroupplot
%             \imagesc{data/spectrogram_noisy.tsv}
% 
%             % Use different color range
%             \nextgroupplot[
%                 title={Different color range},
%                 point meta min=-30,
%                 colormap name=hot2,
%                 point meta max=0]
%             \imagesc{data/spectrogram_noisy.tsv}
%             \node[draw, fill=white, anchor=north west, font=\footnotesize] at (rel axis cs:0, 0.99) {colormap name=hot2};
% 
%             % Show different excerpt
%             \nextgroupplot[
%                 title={Other excerpt},
%                 point meta min=-30,
%                 point meta max=0,
%                 colormap name=parula,
%                 xmin=2.75,
%                 xmax=4.5,
%                 ymin=0,
%                 ymax=8]
%             \imagesc{data/spectrogram_noisy.tsv}
%             \node[draw, fill=white, anchor=north west, font=\footnotesize] at (rel axis cs:0, 0.99) {colormap name=parula};
% 
%             % Use custom defined colormap (it's intentionally ugly)
%             \nextgroupplot[
%                 title={Custom Colormap},
%                 colormap name=temp, 
%                 point meta min=-30,
%                 point meta max=10,
%                 xmin=3.5,
%                 xmax=5.5]
%             \imagesc{data/spectrogram_noisy.tsv}
%             \node[draw, fill=white, anchor=north west, font=\footnotesize] at (rel axis cs:0, 0.99) {colormap name=temp};
%         \end{groupplot}
%     \end{tikzpicture}
% \end{frame}
